% ch03_u02 —— 原书 p.71–p.80（3.3.1 末至 3.5 开头）
\SourcePage{71}

取 $x\in X$。整体 $\mathfrak{g}$-oper 到 $\operatorname{Spec}\mathcal{O}_{x}$ 的限制，定义了仿射概形之间的态射
\[
  \operatorname{Op}_{\mathfrak{g}}(X)\longrightarrow\operatorname{Op}_{\mathfrak{g}}(\mathcal{O}_{x}).
\]
这是闭嵌入，于是得到坐标环之间的满射
\begin{equation}\tag{50}
  \theta_{x}^{A} : A_{\mathfrak{g}}(\mathcal{O}_{x})\twoheadrightarrow A_{\mathfrak{g}}(X).
\end{equation}
$\theta_{x}^{A}$ 与典范滤链严格相容（要看出这一点，可用同构 (24)）。

\bdp{3.3.2} \emph{定理.} 存在唯一的滤代数同构
\begin{equation}\tag{51}
  \varphi_{X} : A_{\mathfrak{g}}(X)\xrightarrow{\ \sim\ }\mathfrak{z}_{{}^{L}\mathfrak{g}}(X),
\end{equation}
使得对任意 $x\in X$，图表
\[
\begin{tikzcd}
  A_{\mathfrak{g}}(\mathcal{O}_{x}) \arrow[r, "\theta_{x}^{A}"] \arrow[d, "\varphi_{\mathcal{O}_{x}}\,\wr"']
    & A_{\mathfrak{g}}(X) \arrow[d, "\wr\,\varphi_{X}"] \\
  \mathfrak{z}_{{}^{L}\mathfrak{g}}(\mathcal{O}_{x}) \arrow[r, "\theta_{x}^{\mathfrak{z}}"']
    & \mathfrak{z}_{{}^{L}\mathfrak{g}}(X)
\end{tikzcd}
\]
交换（这里 $\varphi_{\mathcal{O}_{x}}$ 是 $O=\mathcal{O}_{x}$ 时的同构 (49)）。有 $\sigma_{\mathfrak{z}(X)}\cdot\operatorname{gr}\varphi_{X}=\sigma_{A(X)}$，其中 $\sigma_{A(X)}$ 是同构 (44)，而 $\sigma_{\mathfrak{z}(X)}:\operatorname{gr}\mathfrak{z}(X)\to\mathfrak{z}^{cl}(X)$ 的定义见 2.7.4 末。

\emph{证明.} 由于 $\theta_{x}^{A}$ 与 $\theta_{x}^{\mathfrak{z}}$ 都是满射，且都与滤链严格相容，只需证明存在使图表交换的同构 $\varphi_{X}$。按 2.6.5，有 $\mathcal{D}_{X}$-代数 $\mathcal{A}_{\mathfrak{g}}:=A_{\mathfrak{g}}(\mathcal{O})_{X}$，其纤维为 $A_{\mathfrak{g}}(\mathcal{O}_{x})$。任一整体 oper $\mathfrak{F}\in\operatorname{Op}_{\mathfrak{g}}(X)$ 定义截面 $\gamma_{\mathfrak{F}}:X\to\operatorname{Spec}\mathcal{A}_{\mathfrak{g}}$，其中 $\gamma_{\mathfrak{F}}(x)$ 是 $\mathfrak{F}$ 到 $\operatorname{Spec}\mathcal{O}_{x}$ 的限制。截面 $\gamma_{\mathfrak{F}}$ 是水平的，由此得到 $\operatorname{Op}_{\mathfrak{g}}(X)$ 与 $\operatorname{Spec}\mathcal{A}_{\mathfrak{g}}$ 的水平截面概形之间的同构（认为这需要证明的读者

\SourcePage{72}
可以在 3.3.3 中找到它）。过渡到坐标环，得到典范同构
\begin{equation}\tag{52}
  A_{\mathfrak{g}}(X)\xrightarrow{\ \sim\ }H_{\nabla}(X,\mathcal{A}_{\mathfrak{g}})
\end{equation}
（$H_{\nabla}$ 的定义见 2.6.2）。另一方面，(49) 给出 $\mathcal{D}_{X}$-代数的同构
\[
  \varphi : \mathcal{A}_{\mathfrak{g}}\xrightarrow{\ \sim\ }\mathfrak{z}_{{}^{L}\mathfrak{g}},
\]
从而有同构
\begin{equation}\tag{53}
  H_{\nabla}(X,\mathcal{A}_{\mathfrak{g}})\xrightarrow{\ \sim\ }H_{\nabla}(X,\mathfrak{z}_{{}^{L}\mathfrak{g}})
  =\mathfrak{z}_{{}^{L}\mathfrak{g}}(X).
\end{equation}
现在 $\varphi_{X}$ 就是 (52) 与 (53) 的复合。\hfill$\Box$

\bdp{3.3.3} 本小节（读者完全可以跳过）证明：$\mathfrak{g}$-oper 可以等同于 $\operatorname{Spec}\mathcal{A}_{\mathfrak{g}}$ 的水平截面（3.3.2 用到了这一等同）。

记 $\mathfrak{g}^{+}$ 为所有满足下述条件的 $a\in\mathfrak{g}^{-1}$ 的集合：$a$ 在 $\mathfrak{g}^{\alpha}$ 中的像，对每个单负根 $\alpha$ 都非零（记号见 3.1.1）。$\mathfrak{g}^{+}$ 是仿射概形。考虑 $\operatorname{Aut}^{0}\mathcal{O}$ 通过标准特征标 $\operatorname{Aut}^{0}\mathcal{O}\to\operatorname{Aut}(t\mathcal{O}/t^{2}\mathcal{O})=\mathbb{G}_{m}$ 在 $\mathfrak{g}^{+}$ 上的作用。记 $B$ 为 $\mathfrak{g}$ 所对应的伴随群的 Borel 子群，并给 $B$ 配备 $\operatorname{Aut}^{0}\mathcal{O}$ 的平凡作用。应用 2.6.7 中的函子 $\mathcal{J}:\{\operatorname{Aut}^{0}\mathcal{O}\text{-概形}\}\to\{\operatorname{Aut}\mathcal{O}\text{-概形}\}$，得到：$\mathcal{J}B$ 是态射 $\operatorname{Spec}\mathcal{O}\to B$ 的概形，$\mathcal{J}\mathfrak{g}^{+}$ 是 $\operatorname{Spec}\mathcal{O}$ 上 $\mathfrak{g}^{+}$-值微分形式的概形。群 $\mathcal{J}B$ 通过规范变换作用在 $\mathcal{J}\mathfrak{g}^{+}$ 上，而 $\operatorname{Op}_{\mathfrak{g}}(\mathcal{O})$ 是商概形。$\mathcal{J}B$ 在 $\mathcal{J}\mathfrak{g}^{+}$ 上的作用与态射 $\mathcal{J}\mathfrak{g}^{+}\to\operatorname{Op}_{\mathfrak{g}}(\mathcal{O})$ 都是 $\operatorname{Aut}\mathcal{O}$-等变的。事实上，$\mathcal{J}\mathfrak{g}^{+}$ 是 $\operatorname{Op}_{\mathfrak{g}}(\mathcal{O})$ 上的 $\mathcal{J}B$-torsor（主齐性空间）。此外，选取 $\eta\in\omega_{\mathcal{O}}^{+}:=\omega_{\mathcal{O}}\setminus t\omega_{\mathcal{O}}$ 定义了它的一个截面 $S_{\eta}\subset\mathcal{J}\mathfrak{g}^{+}$：
\[
  S_{\eta}:=\eta\cdot i(f)+V\otimes\omega_{\mathcal{O}}
\]
（这里 $f:=\begin{pmatrix}0&0\\1&0\end{pmatrix}$，而 $i,V$ 已在 3.1.9 中定义）。$S_{\eta}$ 是截面这一事实，正是命题 3.1.10 的局部形式。截面 $S_{\eta}$ 定义了诱导 torsor $\mathfrak{g}^{+}\times\omega_{\mathcal{O}}^{+}\to\operatorname{Op}_{\mathfrak{g}}(\mathcal{O})\times\omega_{\mathcal{O}}^{+}$ 的 $\operatorname{Aut}\mathcal{O}$-等变截面
\[
  s : \operatorname{Op}_{\mathfrak{g}}(\mathcal{O})\times\omega_{\mathcal{O}}^{+}\longrightarrow\mathfrak{g}^{+}\times\omega_{\mathcal{O}}^{+}.
\]

\SourcePage{73}
现在考虑 $\mathcal{D}_{X}$-概形 $(\mathcal{J}\mathfrak{g}^{+})_{X}$、$(\mathcal{J}B)_{X}$ 与 $\operatorname{Op}_{\mathfrak{g}}(\mathcal{O})_{X}=\operatorname{Spec}\mathcal{A}_{\mathfrak{g}}$。显然 $(\mathcal{J}B)_{X}$ 是 $X$ 上的群 $\mathcal{D}_{X}$-概形，而概形 $(\mathcal{J}\mathfrak{g}^{+})_{X}$ 是 $\operatorname{Op}_{\mathfrak{g}}(\mathcal{O})_{X}$ 上的 $(\mathcal{J}B)_{X}$-torsor。事实上 $(\mathcal{J}B)_{X}=\mathcal{J}(B_{X})$，而 $(\mathcal{J}\mathfrak{g}^{+})_{X}$ 是 $X$ 上 $\mathfrak{g}^{+}$-值微分形式的 jets 概形。显然
\[
  \operatorname{Op}_{\mathfrak{g}}=\operatorname{Sect}(\mathfrak{g}_{X}^{+})/\operatorname{Sect}(B_{X})
  =\operatorname{Sect}^{\nabla}((\mathcal{J}\mathfrak{g}^{+})_{X})/\operatorname{Sect}^{\nabla}((\mathcal{J}B)_{X})
  \subset\operatorname{Sect}^{\nabla}(\operatorname{Op}_{\mathfrak{g}}(\mathcal{O})_{X}).
\]
这里 Sect 表示 $X$-概形的截面层，$\operatorname{Sect}^{\nabla}$ 表示 $\mathcal{D}_{X}$-概形的水平截面层。为证明 $\operatorname{Op}_{\mathfrak{g}}=\operatorname{Sect}^{\nabla}(\operatorname{Op}_{\mathfrak{g}}(\mathcal{O})_{X})$，只需证明 $\operatorname{Sect}^{\nabla}((\mathcal{J}\mathfrak{g}^{+})_{X})\to\operatorname{Sect}^{\nabla}(\operatorname{Op}_{\mathfrak{g}}(\mathcal{O})_{X})$ 是满射。为此，使用由 $s$ 诱导的 $\mathcal{D}_{X}$-概形态射 $\operatorname{Op}_{\mathfrak{g}}(\mathcal{O})_{X}\times(\omega_{\mathcal{O}}^{+})_{X}\to(\mathfrak{g}^{+})_{X}$，以及如下事实：$(\omega_{\mathcal{O}}^{+})_{X}$（即处处非零的微分形式的 jets 概形）在 $X$ 的每一点的邻域内都有水平截面。

于是我们把 $\operatorname{Op}_{\mathfrak{g}}(X)$ 等同于 $\operatorname{Op}_{\mathfrak{g}}(\mathcal{O})_{X}=\operatorname{Spec}\mathcal{A}_{\mathfrak{g}}$ 的水平截面的集合。同样地，可把概形 $\operatorname{Op}_{\mathfrak{g}}(X)$ 等同于 $\operatorname{Spec}\mathcal{A}_{\mathfrak{g}}$ 的水平截面概形。

\emph{注.} 我们使用 $s$ 只是为了简化证明 ???。
% [Tutor Note: 原书 p.73 此处引用印作“???”，底本如此，保留原样。]

\bdsec{3.4. $G$-opers 与 $\mathfrak{g}$-opers}
本小节假设 $G$ 半单（其实一般情形也可类似处理，见 3.4.2 末的注 iii）。对每个单负根 $\alpha$，固定非零元 $y_{\alpha}\in\mathfrak{g}^{\alpha}$。令 $G_{ad}:=G/Z$，$B_{ad}:=B/Z$，$H_{ad}:=H/Z$，其中 $Z$ 是 $G$ 的中心。

\bdp{3.4.1} 存在显然的投影 $\operatorname{Op}_{G}(X)\longrightarrow\operatorname{Op}_{\mathfrak{g}}(X):=\operatorname{Op}_{G_{ad}}(X)$。我们将构造一个截面 $\operatorname{Op}_{\mathfrak{g}}(X)\to\operatorname{Op}_{G}(X)$，它依赖于 $\omega_{X}$ 的平方根的选取，即依赖于带同构 $\mathcal{L}^{\otimes 2}\xrightarrow{\sim}\omega_{X}$ 的线丛 $\mathcal{L}$。设 $(\mathfrak{F}_{B_{ad}},\nabla)$ 是 $\mathfrak{g}$-oper。把它提升为 $G$-oper，等价于把 $\mathfrak{F}_{B_{ad}}$ 提升为 $B$-丛，这又等价于把 $\mathfrak{F}_{H_{ad}}$ 提升为 $H$-丛（这里 $\mathfrak{F}_{H_{ad}}$ 是 $\mathfrak{F}_{B_{ad}}$ 沿 $B_{ad}\to H_{ad}$ 的前推）。在 $\mathfrak{g}=\mathfrak{sl}_{2}$ 的特殊情形，我们在 3.1.8 中构造了典范同构 $\mathfrak{F}_{H_{ad}}\xrightarrow{\sim}\omega_{X}$；3.1.8 的构造用到了固定元 $f\in\mathfrak{sl}_{2}/\mathfrak{b}_{0}$。完全类似地，

\SourcePage{74}
一般情形下可构造典范同构 $\mathfrak{F}_{H_{ad}}\xrightarrow{\sim}\lambda\omega_{X}$，其中 $\lambda\omega_{X}$ 是 $\omega_{X}$ 沿同态 $\lambda:\mathbb{G}_{m}\to H_{ad}$ 的前推，而 $\lambda$ 使得对任意单正根 $\alpha$，$\lambda(t)$ 在 $\mathfrak{g}^{\alpha}$ 上乘以 $t$（构造用到 3.4 开头固定的那些元素 $y_{\alpha}$）。存在唯一的同态 $\lambda^{\#}:\mathbb{G}_{m}\to H$，使得
\begin{equation}\tag{54}
  \lambda^{\#}(t)\ \operatorname{mod}Z=\lambda(t)^{2}
\end{equation}
（事实上，$\lambda$ 对应余权 $\check{\rho}:=$ 基本余权之和，而 $2\check{\rho}$ 属于余根格）。我们把 $\mathfrak{F}_{H_{ad}}=\lambda\omega_{X}$ 提升为 $H$-丛 $\lambda^{\#}\mathcal{L}$，其中 $\mathcal{L}$ 是 $\omega_{X}$ 的平方根。

\bdp{3.4.2} 记 $\omega^{1/2}(X)$ 为 $\omega_{X}$ 的平方根的群胚。对固定的 $\mathcal{L}\in\omega^{1/2}(X)$，有等价
\begin{equation}\tag{55}
  \Phi_{\mathcal{L}} : \operatorname{Op}_{\mathfrak{g}}(X)\times Z\operatorname{tors}(X)\xrightarrow{\ \sim\ }\operatorname{Op}_{G}(X),
\end{equation}
其中 $Z\operatorname{tors}(X)$ 是 $X$ 上 $Z$-torsor 的群胚。$\Phi_{\mathcal{L}}(\mathfrak{F},\mathcal{E})$ 的定义如下：用 $\mathcal{L}$ 把 $\mathfrak{F}\in\operatorname{Op}_{\mathfrak{g}}(X)$ 提升为 $G$-oper（见 3.4.1），再用 $\mathcal{E}$ 扭这个 $G$-oper。$\Phi_{\mathcal{L}}$ 对 $\mathcal{L}$ 的依赖方式为
\[
  \Phi_{\mathcal{L}\otimes\mathcal{A}}(\mathfrak{F},\mathcal{E})=\Phi_{\mathcal{L}}(\mathfrak{F},\mathcal{E}\cdot\alpha\mathcal{A}).
\]
这里 $\mathcal{A}$ 是 $\mathcal{O}_{X}$ 的平方根，或者说，是 $X$ 上的 $\mu_{2}$-torsor；而 $\alpha\mathcal{A}$ 是 $\mu_{2}$-torsor $\mathcal{A}$ 沿态射
\begin{equation}\tag{56}
  \alpha : \mu_{2}\longrightarrow Z,\qquad \alpha:=\lambda^{\#}|_{\mu_{2}}
\end{equation}
的前推。回忆 $\lambda^{\#}$ 由 (54) 定义。

\emph{注.}
\begin{enumerate}[label=(\roman*),leftmargin=2.2em]
  \item 若把 $\operatorname{Op}_{\mathfrak{g}}(X)$ 看作概形，把 $\operatorname{Op}_{G}(X)$ 与 $Z\operatorname{tors}(X)$ 看作代数叠，则 (55) 成为代数叠的同构。
  \item $\alpha$ 是“那个”主同态 $\mathrm{SL}_{2}\to G$ 到中心 $\mu_{2}\subset\mathrm{SL}_{2}$ 的限制。

\SourcePage{75}
  \item 若 $G$ 约化但非半单，而 $\mathfrak{g}:=\operatorname{Lie}(G/Z)$，则定义依赖于 $\mathcal{L}\in\omega^{1/2}(X)$ 的截面 $\operatorname{Op}_{\mathfrak{g}}(X)\to\operatorname{Op}_{G}(X)$ 为复合 $\operatorname{Op}_{\mathfrak{g}}(X)\to\operatorname{Op}_{[G,G]}(X)\to\operatorname{Op}_{G}(X)$。如果把 $Z\operatorname{tors}(X)$ 换成 $Z^{\nabla}\operatorname{tors}(X)$，即 $X$ 上带联络的 $Z$-torsor 的群胚，则 3.4.2 的结果仍然成立。
\end{enumerate}

\bdp{3.4.3} 下面是 3.4.2 的一个更自然的表述。先引入群胚 $Z\operatorname{tors}_{\theta}(X)$（$\theta$ 应使读者想到 $\theta$-特征，即 $\omega_{X}$ 的平方根）。$Z\operatorname{tors}_{\theta}(\mathcal{L})$ 的对象是偶对 $(\mathcal{E},\mathcal{L})$，其中 $\mathcal{E}\in Z\operatorname{tors}(X)$，$\mathcal{L}\in\omega^{1/2}(X)$，但我们宁愿写 $\mathcal{E}\cdot\mathcal{L}$ 而不写 $(\mathcal{E},\mathcal{L})$。
% [Tutor Note: 原书 p.75 此处印作 Ztors_{\theta}(\mathcal{L})，按上下文应为 Ztors_{\theta}(X)，底本如此，保留原样。]
令
\[
  \operatorname{Mor}(\mathcal{E}_{1}\cdot\mathcal{L}_{1},\mathcal{E}_{2}\cdot\mathcal{L}_{2})
  :=\operatorname{Mor}(\mathcal{E}_{1},\mathcal{E}_{2}\cdot\alpha(\mathcal{L}_{2}/\mathcal{L}_{1})),
\]
其中 $\alpha(\mathcal{L}_{2}/\mathcal{L}_{1})$ 是 $\mu_{2}$-torsor $\mathcal{L}_{2}/\mathcal{L}_{1}:=\mathcal{L}_{2}\otimes\mathcal{L}_{1}^{\otimes(-1)}$ 沿同态 (56) 的前推。态射的复合按显然方式定义。于是可把 3.4.2 重述为典范等价
\begin{equation}\tag{57}
  \Phi : \operatorname{Op}_{\mathfrak{g}}(X)\times Z\operatorname{tors}_{\theta}(X)\xrightarrow{\ \sim\ }\operatorname{Op}_{G}(X),
\end{equation}
其中 $\Phi(\mathfrak{F},\mathcal{L}\cdot\mathcal{E}):=\Phi_{\mathcal{L}}(\mathfrak{F},\mathcal{E})$，而 $\Phi_{\mathcal{L}}$ 由 (55) 定义。

在 3.2.1 的局部情形有类似的典范等价
\begin{equation}\tag{58}
  \operatorname{Op}_{\mathfrak{g}}(\mathcal{O})\times Z\operatorname{tors}_{\theta}(\mathcal{O})\xrightarrow{\ \sim\ }\operatorname{Op}_{G}(\mathcal{O}),
\end{equation}
其中 $Z\operatorname{tors}_{\theta}(\mathcal{O})$ 按整体情形定义。当然，$Z\operatorname{tors}_{\theta}(\mathcal{O})$ 的所有对象彼此同构，且一个对象的自同构群是 $Z$。对 $Z\operatorname{tors}(\mathcal{O})$ 同样如此。若把 $\mathcal{O}$ 的自同构也考虑进来，$Z\operatorname{tors}_{\theta}(\mathcal{O})$ 与 $Z\operatorname{tors}(\mathcal{O})$ 的差别就清楚了（见 3.5.2）。

\bdp{3.4.4} 要描述 $Z\operatorname{tors}_{\theta}(\mathcal{O})$ 的“经济”版本，我们需要一些抽象废话。

设 $Z$ 是交换群。范畴 $C$ 上的 $Z$-结构是态射 $Z\to\operatorname{Aut}\mathrm{id}_{C}$。等价地，$C$ 上的 $Z$-结构是 $Z$ 在 $\operatorname{Mor}(c_{1},c_{2})$（$c_{1},c_{2}\in\operatorname{Ob}C$）上的作用，使得对任意态射 $c_{1}\xrightarrow{f}c_{2}\xrightarrow{g}c_{3}$ 与任意 $z\in Z$，有 $z(gf)=(zg)f=g(zf)$。$Z$-范畴是带

\SourcePage{76}
$Z$-结构的范畴。设 $C$ 与 $C'$ 是 $Z$-范畴。若对任意 $c_{1},c_{2}\in C$，映射 $\operatorname{Mor}(c_{1},c_{2})\to\operatorname{Mor}(F(c_{1}),F(c_{2}))$ 是 $Z$-等变的，则称函子 $F:C\to C'$ 为 $Z$-函子。若 $A\to Z$ 是交换群的态射，且 $C$ 是 $A$-范畴，则如下定义诱导的 $Z$-范畴 $C\otimes_{A}Z$：$\operatorname{Ob}(C\otimes_{A}Z)=\operatorname{Ob}C$，$(C\otimes_{A}Z)$-态射 $c_{1}\to c_{2}$ 的集合为
\[
  (\operatorname{Mor}_{C}(c_{1},c_{2})\times Z)/A
  =\{\text{由 }A\text{-集 }\operatorname{Mor}_{C}(c_{1},c_{2})\text{ 诱导的 }Z\text{-集}\},
\]
而 $C\otimes_{A}Z$ 中态射的复合按显然方式定义。有自然的 $A$-函子 $C\to C\otimes_{A}Z$；并且对任意 $Z$-范畴 $C'$，任一 $A$-函子 $C\to C'$ 有唯一分解 $C\to C\otimes_{A}Z\xrightarrow{F}C'$，其中 $F$ 是 $Z$-函子。

记 $\omega^{1/2}(\mathcal{O})$ 为 $\omega_{\mathcal{O}}$ 的平方根的群胚。这是 $\mu_{2}$-范畴。$Z\operatorname{tors}(\mathcal{O})$ 与 $Z\operatorname{tors}_{\theta}(\mathcal{O})$ 是 $Z$-范畴。典范 $\mu_{2}$-函子 $\omega^{1/2}(\mathcal{O})\to Z\operatorname{tors}_{\theta}(\mathcal{O})$ 诱导等价
\[
  \omega^{1/2}(\mathcal{O})\otimes_{\mu_{2}}Z\to Z\operatorname{tors}_{\theta}(\mathcal{O}).
\]

\bdp{3.4.5} 读者也许更喜欢 $Z\operatorname{tors}_{\theta}(X)$ 与 $Z\operatorname{tors}_{\theta}(\mathcal{O})$ 的如下“具体”版本。定义正合列
\begin{equation}\tag{59}
  0\to Z\to\widetilde{Z}\to\mathbb{G}_{m}\to 0
\end{equation}
为
\begin{equation}\tag{60}
  0\to\mu_{2}\to\mathbb{G}_{m}\xrightarrow{f}\mathbb{G}_{m}\to 0,\qquad f(x):=x^{2}
\end{equation}
沿态射 (56) 的前推。记 $\widetilde{Z}\operatorname{tors}_{\omega}(X)$ 为把 $\mathbb{G}_{m}$-torsor $\omega_{X}$ 提升为 $\widetilde{Z}$-torsor 所成的群胚（即 $\widetilde{Z}\operatorname{tors}_{\omega}(X)$ 的对象，是 $X$ 上的 $\widetilde{Z}$-torsor，连同相应的 $\mathbb{G}_{m}$-torsor 与 $\omega_{X}$ 之间的一个同构）。由 (60) 到 (59) 的态射诱导函子 $F:\omega^{1/2}(X)\to\widetilde{Z}\operatorname{tors}_{\omega}(X)$。由
\[
  \mathcal{E}\cdot\mathcal{L}\mapsto\mathcal{E}\cdot F(\mathcal{L}),\qquad
  \mathcal{E}\in Z\operatorname{tors}(X),\ \mathcal{L}\in\omega^{1/2}(X)
\]
定义的函子 $Z\operatorname{tors}_{\theta}(X)\to\widetilde{Z}\operatorname{tors}_{\omega}(X)$ 是等价。

完全类似地可定义 $\widetilde{Z}\operatorname{tors}_{\omega}(\mathcal{O})$ 与典范等价 $Z\operatorname{tors}_{\theta}(\mathcal{O})\xrightarrow{\sim}\widetilde{Z}\operatorname{tors}_{\omega}(\mathcal{O})$。

\SourcePage{77}
借助 $\widetilde{Z}\operatorname{tors}_{\omega}(X)$ 与 $\widetilde{Z}\operatorname{tors}_{\omega}(\mathcal{O})$，可以容易地理解等价 (57) 与 (58)。例如，构造等价
\[
  \operatorname{Op}_{\mathfrak{g}}(X)\times\widetilde{Z}\operatorname{tors}_{\omega}(X)\xrightarrow{\ \sim\ }\operatorname{Op}_{G}(X).
\]
考虑如下行正合的交换图表：
\begin{equation}\tag{61}
  \begin{tikzcd}
    0\arrow[r] & \mu_{2}\arrow[r]\arrow[d, "\alpha"'] & \mathbb{G}_{m}\arrow[r]\arrow[d, "\lambda^{\#}"'] & \mathbb{G}_{m}\arrow[r]\arrow[d, "\lambda"] & 0\\
    0\arrow[r] & Z\arrow[r] & H\arrow[r] & H_{ad}\arrow[r] & 0
  \end{tikzcd}
\end{equation}
这里上行是 (60)；下行与态射 $\lambda,\lambda^{\#}$ 在 3.4.1 中定义。按 3.4.1，$X$ 上的 $G$-oper 等同于 $X$ 上的 $\mathfrak{g}$-oper 加上把 $H_{ad}$-torsor $\lambda_{*}(\omega_{X})$ 提升为 $H$-torsor。这样的提升等同于 $\widetilde{Z}\operatorname{tors}_{\omega}(X)$ 的一个对象：考察交换图表
\begin{equation}\tag{62}
  \begin{tikzcd}
    0\arrow[r] & Z\arrow[r]\arrow[d, equal] & \widetilde{Z}\arrow[r]\arrow[d] & \mathbb{G}_{m}\arrow[r]\arrow[d, "\lambda"] & 0\\
    0\arrow[r] & Z\arrow[r] & H\arrow[r] & H_{ad}\arrow[r] & 0
  \end{tikzcd}
\end{equation}
的右端（笛卡尔）方块（(62) 的上行是 (59)，而 (62) 与 (61) 的下行相同）。

\bdp{3.4.6} $Z\operatorname{tors}(X)$ 是（严格交换的）Picard 范畴（见 [Del73] 中的定义 1.4.2），而 $Z\operatorname{tors}_{\theta}(X)$ 是 $Z\operatorname{tors}(X)$ 上的一个 torsor；确切地说，$Z\operatorname{tors}_{\theta}(X)$ 是由 $\mu_{2}\operatorname{tors}(X)$ 上的 torsor $\omega^{1/2}(X)$ 沿对应于 (56) 的 Picard 函子 $\mu_{2}\operatorname{tors}(X)\to Z\operatorname{tors}(X)$ 诱导而来。我们在 §4 中将使用这套语言，所以这里回忆一下定义。

\emph{Picard 范畴}是指 [De-Mi] 意义下的张量范畴 $\mathcal{A}$（即对称交换幺半范畴），使得 $\mathcal{A}$ 的所有态射可逆（即 $\mathcal{A}$ 是群胚），且 $\mathcal{A}$ 的所有对象可逆；也就是说，对每个 $a\in\operatorname{Ob}\mathcal{A}$，存在 $a'\in\operatorname{Ob}\mathcal{A}$，使 $a\cdot a'$ 是单位对象（我们用 $\cdot$ 记“张量积”函子 $\mathcal{A}\times\mathcal{A}\to\mathcal{A}$；在 [De-Mi] 与 [Del73] 中它分别记作 $\otimes$ 与 $+$）。严格交换是指

\SourcePage{78}
对每个 $a\in\operatorname{Ob}\mathcal{A}$，交换同构 $a\otimes a\xrightarrow{\sim}a\otimes a$ 是恒等。

幺半范畴 $\mathcal{A}$ 在范畴 $C$ 上的一个\emph{作用}是函子 $\mathcal{A}\times C\to C$（记作 $\cdot$），连同结合约束，即函子同构 $(a_{1}\cdot a_{2})\cdot c\xrightarrow{\sim}a_{1}\cdot(a_{2}\cdot c)$（$a_{i}\in\mathcal{A}$，$c\in C$），它满足一个五边形公理，该公理与 $\mathcal{A}$ 中结合约束的五边形公理类似（见 [Del73] 与 [De-Mi]）；我们还要求对应于 $\mathcal{A}$ 的单位对象的函子 $F:C\to C$ 全忠实（此时同构 $F^{2}\xrightarrow{\sim}F$ 给出典范同构 $F\xrightarrow{\sim}\mathrm{id}$）。这一定义可在 [Pa] 与 [Yet] 的 §3 中找到。\emph{$\mathcal{A}$-模}是带 $\mathcal{A}$ 的作用的范畴。若 $C$ 与 $\widetilde{C}$ 是 $\mathcal{A}$-模，则\emph{$\mathcal{A}$-模函子} $C\to\widetilde{C}$ 是带函子同构 $\Phi(a\cdot c)\xrightarrow{\sim}a\cdot\Phi(c)$ 的函子 $\Phi:C\to\widetilde{C}$，且满足自然的相容条件（构造同构 $\Phi((a_{1}\cdot a_{2})\cdot c)\xrightarrow{\sim}a_{1}\cdot(a_{2}\cdot\Phi(c))$ 的两种方式必须给出同一结果；见 [Pa]、[Yet]）。$\mathcal{A}$-模函子也称为 $\mathcal{A}$-模的\emph{态射}。

Picard 范畴 $\mathcal{A}$ 上的一个\emph{torsor} 是 $\mathcal{A}$-模，使得对某个 $c\in\operatorname{Ob}C$，函子 $a\mapsto a\cdot c$ 是 $\mathcal{A}$ 与 $C$ 之间的等价（此时对一切 $c\in\operatorname{Ob}C$ 都成立）。

设 $\mathcal{A}$ 与 $\mathcal{B}$ 是 Picard 范畴。\emph{Picard 函子} $\mathcal{A}\to\mathcal{B}$ 是带函子同构 $F(a_{1}\cdot a_{2})\xrightarrow{\sim}F(a_{1})\cdot F(a_{2})$ 的函子 $F:\mathcal{A}\to\mathcal{B}$，且该同构与交换约束、结合约束相容。于是 $F$ 把 $\mathcal{A}$ 的单位对象送到 $\mathcal{B}$ 的单位对象，即 $F$ 是 [De-Mi] 意义下的张量函子。在 [Del73] 中，Picard 函子称为加法函子。

设 $F:\mathcal{A}_{1}\to\mathcal{A}_{2}$ 是 Picard 函子，$C_{i}$ 是 $\mathcal{A}_{i}$ 上的 torsor（$i=1,2$）。则 $C_{2}$ 带有 $\mathcal{A}_{1}$ 的作用。在这种情形下，$\mathcal{A}_{1}$-模函子 $C_{1}\to C_{2}$ 称为 $F$-\emph{仿射}函子。

\emph{例.}
\begin{enumerate}[label=\arabic*),leftmargin=2.2em]
  \item 设 $A$ 是交换代数群。则 $A\operatorname{tors}(X)$ 有典范的 Picard 范畴结构。

\SourcePage{79}
  \item 交换代数群的态射 $A\to B$ 诱导 Picard 函子 $A\operatorname{tors}(X)\to B\operatorname{tors}(X)$。
  \item 3.4.2 中的群胚 $\omega^{1/2}(X)$ 是 Picard 范畴 $\mu_{2}\operatorname{tors}(X)$ 上的 torsor。3.4.3 中的群胚 $Z\operatorname{tors}_{\theta}(X)$ 是 $Z\operatorname{tors}(X)$ 上的 torsor。
\end{enumerate}

若 $F:\mathcal{A}\to\mathcal{B}$ 是 Picard 范畴之间的 Picard 函子，$C$ 是 $\mathcal{A}$ 上的 torsor，则可得 $\mathcal{B}$ 上的\emph{诱导} torsor $\mathcal{B}\cdot_{\mathcal{A}}C$。读者可以从下例复原 $\mathcal{B}\cdot_{\mathcal{A}}C$ 的定义（见 3.4.3）：若 $\mathcal{A}=\mu_{2}\operatorname{tors}(X)$，$\mathcal{B}=Z\operatorname{tors}(X)$，$F$ 来自 (56)，$C=\omega^{1/2}(X)$，则 $\mathcal{B}\cdot_{\mathcal{A}}C=Z\operatorname{tors}_{\theta}(X)$。$\mathcal{B}\cdot_{\mathcal{A}}C$ 的对象记作 $b\cdot c$（$b\in\operatorname{Ob}\mathcal{B}$，$c\in\operatorname{Ob}C$，见 3.4.3）。

感兴趣的读者可以表述 $\mathcal{B}\cdot_{\mathcal{A}}C$ 的泛性质。我们需要如下较弱的性质。给定带 $\mathcal{B}$ 作用的范畴 $\widetilde{C}$ 与 $\mathcal{A}$-模函子 $\Phi:C\to\widetilde{C}$，有自然的方式构造 $\mathcal{B}$-模函子 $\Psi:\mathcal{B}\cdot_{\mathcal{A}}C\to\widetilde{C}$：令 $\Psi(b\cdot c):=b\cdot\Phi(c)$，并在态射上按显然方式定义 $\Psi$（即映射 $\operatorname{Mor}(b_{1}\cdot c_{1},b_{2}\cdot c_{2})\to\operatorname{Mor}(b_{1}\cdot\Phi(c_{1}),b_{2}\cdot\Phi(c_{2}))$ 是复合
\[
  \operatorname{Mor}(b_{1}\cdot c_{1},b_{2}\cdot c_{2})
  =\operatorname{Mor}(b_{1},b_{2}\cdot c_{2}/c_{1})
  \to\operatorname{Mor}(b_{1}\cdot\Phi(c_{1}),b_{2}\cdot c_{2}/c_{1}\cdot\Phi(c_{1}))
  \xrightarrow{\sim}\operatorname{Mor}(b_{1}\cdot\Phi(c_{1}),b_{2}\cdot\Phi(c_{2}))).
\]
同构 $\Psi(b_{1}\cdot(b_{2}\cdot c))\xrightarrow{\sim}b_{1}\cdot\Psi(b_{2}\cdot c)$ 是显然的那一个。

我们在下述情形使用这一构造。设有 Picard 函子 $\ell:\mathcal{B}\to\widetilde{\mathcal{B}}$、$\widetilde{\mathcal{B}}$ 上的 torsor $\widetilde{C}$、以及 $\ell'$-仿射函子 $\Phi:C\to\widetilde{C}$，其中 $\ell'$ 是复合 $\mathcal{A}\xrightarrow{F}\mathcal{B}\xrightarrow{\ell}\widetilde{\mathcal{B}}$。则上述构造给出 $\ell$-仿射函子 $\mathcal{B}\cdot_{\mathcal{A}}C\to\widetilde{C}$。

\bdp{3.4.7} 设 $Z$ 是交换群，$Z\operatorname{tors}$ 是 $Z$-torsor 的 Picard 范畴（一点上的）。下述注记将在 4.4.9 中使用。

\emph{注.}
\begin{enumerate}[label=(\roman*),leftmargin=2.2em]
  \item 从 $Z\operatorname{tors}$ 到 Picard 范畴 $\mathcal{A}$ 的 Picard 函子，“等同于”一个态射 $Z\to\operatorname{Aut}\mathbf{1}_{\mathcal{A}}$，其中 $\mathbf{1}_{\mathcal{A}}$ 是 $\mathcal{A}$ 的单位对象。更确切地说，从 Picard 函子范畴

\SourcePage{80}
$Z\operatorname{tors}\to\mathcal{A}$ 到 $\operatorname{Hom}(Z,\operatorname{Aut}\mathbf{1}_{\mathcal{A}})$ 的自然函子是等价。这里集合 $\operatorname{Hom}(Z,\operatorname{Aut}\mathbf{1}_{\mathcal{A}})$ 视为离散范畴。
  \item 把“Picard”换成“幺半”后，前述注记仍然成立。
  \item $Z\operatorname{tors}$ 在范畴 $C$ 上的一个作用，“等同于” $C$ 上的 $Z$-结构，即态射 $Z\to\operatorname{Aut}\mathrm{id}_{C}$（注意幺半范畴 $\mathcal{A}$ 在 $C$ 上的作用等同于幺半函子 $\mathcal{A}\to\operatorname{Funct}(C,C)$，再应用前一注记）。
  \item 设 $C_{1}$ 与 $C_{2}$ 是 $Z\operatorname{tors}$ 上的模。按前一注记，$C_{1}$ 与 $C_{2}$ 是 3.4.4 意义下的 $Z$-范畴。容易看出，$(Z\operatorname{tors})$-模函子 $C_{1}\to C_{2}$ 与 3.4.4 意义下的 $Z$-函子相同（即函子 $F:C_{1}\to C_{2}$ 至多有一个 $(Z\operatorname{tors})$-模函子结构，且这样的结构存在当且仅当 $F$ 是 $Z$-函子）。
  \item $Z\operatorname{tors}$ 上的 torsor “等同于”一个与 $Z\operatorname{tors}$ $Z$-等价的 $Z$-范畴。（我们需要这一条吗？？？）
% [Tutor Note: 原书 p.80 括注“(do we need this???)”为作者存疑之语，底本如此，保留原样。]
\end{enumerate}

\bdsec{3.5. 局部 opers II}
对大多数李代数 $\mathfrak{g}$（例如 $\mathfrak{g}=\mathfrak{sl}_{n}$，$n>5$），定理 3.2.2 中提到的性质并不能\emph{唯一地}确定 Feigin–Frenkel 同构 (49)。原因是 $A_{\mathfrak{g}}(\mathcal{O})$ 有许多在 $\operatorname{gr}A_{\mathfrak{g}}(\mathcal{O})$ 上诱导恒等的 $\operatorname{Aut}\mathcal{O}$-等变自同构。这一点从 3.2.1 中 $\operatorname{Op}_{\mathfrak{g}}(\mathcal{O})=\operatorname{Spec}A_{\mathfrak{g}}(\mathcal{O})$ 的几何描述即可看出，也可从将在 3.5.6 中给出的 $A_{\mathfrak{g}}(\mathcal{O})$ 的描述看出（见 (65)–(68)）。3.5–3.6 的目标是表述唯一确定 Feigin–Frenkel 同构的性质 3.6.7。这一性质以及 3.6.11 将用于我们主定理 5.2.6 的证明。在 3.7 与 3.8 中我们说明如何从 [FF92] 提取性质 3.6.7 与 3.6.11。人们可以（或许应该）在 3.5–3.8 之前先读 §4 与（§5 的一大部分？）。我们当然建议读者，在 ?? 中用到 3.6.11 之前，先跳过 3.5.16–3.5.18 与 3.6.8–3.6.11。
% [Tutor Note: 原书 p.80 有作者留待补全的“?”与残缺引用“??”，底本如此，保留原样。]
